\begin{proof}
Let
\[
 v_{r,k}=\sqrt q\sum_{n\ge0}r^n e_{a^n b^k},\qquad
 u_{r,0}=\sqrt q\sum_{n\ge0}r^n e_{a^n}.
\]
The target's singular decomposition gives
$\Bnorm{v_{r,k}}=\Bnorm{u_{r,0}}=1$ and
$u_{r,0}^*H_{f_{r,k}}=v_{r,k}^*$. If $P_G$ projects onto
$\mathcal S_G$, then $G(I-P_G)=0$, and therefore
\begin{equation}
 \Bnorm{H_{f_{r,k}}-G}
   \ge\sqrt{1-\Bnorm{P_Gv_{r,k}}^2}.
 \label{I:B:eq:boundary-kernel-lower}
\end{equation}
The coordinates of $x\in\mathcal S_G$ satisfy
$x[a^n b^k]=\ell B^kA^nx$. Hence
\begin{equation}
 \Bnorm{P_Gv_{r,k}}
 =\Bnorm{\sqrt q\,\ell B^k(I-rA)^{-1}}
 \le\Bnorm{\ell B^{k-1}},
 \label{I:B:eq:boundary-overlap-bound}
\end{equation}
where the final inequality is \eqref{I:B:eq:resolvent-defect-contraction}.

Consider any sequence $r_j\uparrow1$ and arbitrary rank-at-most-$k$
competitors $G_j$. Zero-rank terms have zero overlap. For all remaining
terms, choose orthonormal coordinates on their row spaces and pad by
zero coordinates to dimension $k$. The padded tuples still satisfy
\eqref{I:B:eq:boundary-passivity}. Passing to a subsequence gives
$A_j\to A_0$, $B_j\to B_0$, $\ell_j\to\ell_0$.

If $I-A_0$ is invertible, the resolvent in
\eqref{I:B:eq:boundary-overlap-bound}'s middle expression stays bounded,
so that overlap tends to zero. Otherwise
$\mathcal E=\ker(I-A_0)$ is nonzero and reduces the contraction $A_0$.
Indeed $A_0x=x$ and contraction imply $A_0^*x=x$.
The limiting passive inequality forces $B_0\mathcal E=0$ and
$\ell_0\mathcal E=0$. In the orthogonal decomposition
$\mathcal E\oplus\mathcal F$ one can consequently write
\[
 B_0=\begin{pmatrix}0&C\\0&D\end{pmatrix},\qquad
 \ell_0=(0,l),\qquad \dim\mathcal F\le k-1.
\]
The $\mathcal F$-column part of passivity gives
$D^*D+l^*l\preceq I_{\mathcal F}$. If $k\ge2$,
$\ell_0B_0^{k-1}=(0,lD^{k-1})$ has norm at most $\Gamma_{k-1}$.
If $k=1$, $\mathcal F=\{0\}$ and $\ell_0=0$, giving the same conclusion
with $\Gamma_0=0$.

Thus every subsequential limiting overlap is at most $\Gamma_{k-1}$.
A compactness contradiction makes this uniform over all competitors:
\[
 \limsup_{r\uparrow1}\sup_{\substack{G\text{ bounded word-Hankel}\\
                                  \operatorname{rank}G\le k}}
       \Bnorm{P_Gv_{r,k}}\le\Gamma_{k-1}.
\]
Insert this bound into \eqref{I:B:eq:boundary-kernel-lower} and use
Lemma~\ref{I:B:lem:one-letter-projection}. No existence assumption for a
best word-Hankel competitor is required.
\end{proof}
